\chapter{Summer}\label{app:appendix01}

Dette bilag præsenterer metoder til at beregne summer, som ofte optræder i analysen af algoritmer. Mange
af formlerne her findes i enhver lærebog i infinitesimalregning, men det er praktisk at have metoderne samlet ét sted. Indholdet i bilaget bygger på \textit{Introduction to Algorithms} \cite{clrs4}. Bilaget afsluttes med et nyttigt oversigtsark til tidskompleksitetsanalyse.

Når en algoritme indeholder en iterativ kontrolstruktur som en løkke af typen \While eller \For\unskip, kan du udtrykke dens køretid som summen af den tid, der bruges på hver udførelse af løkkens krop. I \autoref{sec:loop-analysis} argumenterede vi for eksempel for, at den $i$'te iteration af $\proc{Quadratic-Sum}$ i værste tilfælde tog tid proportional med $i$. Ved at lægge den tid sammen, der blev brugt på hver iteration, fik vi summen (eller rækken) $\sum\limits_{i=1}^n i$. Beregningen af denne sum gav en grænse på $\mathcal{O}(n^2)$ for algoritmens køretid i værste tilfælde. Eksemplet viser, hvorfor du bør kunne regne med summer. Før vi går i detaljer, giver vi et overblik over de vigtigste summer og deres lukkede former.

\begin{table}[htbp]
    \centering
    \renewcommand{\arraystretch}{2}
    \begin{tabular}{|c|c|}
    \hline
    \textbf{Sum} & \textbf{Lukket form} \\ \hline
    $\displaystyle \sum_{k=0}^n ar^k \, (r \neq 1)$ & $\displaystyle \frac{ar^{n+1} - a}{r - 1}$ \\ \hline
    $\displaystyle \sum_{k=1}^n k$ & $\displaystyle \frac{n(n+1)}{2}$ \\ \hline
    $\displaystyle \sum_{k=1}^n k^2$ & $\displaystyle \frac{n(n+1)(2n+1)}{6}$ \\ \hline
    $\displaystyle \sum_{k=1}^n k^3$ & $\displaystyle \frac{n^2(n+1)^2}{4}$ \\ \hline
    $\displaystyle \sum_{k=0}^n x^k \, (x \neq 1)$ & $\displaystyle \frac{x^{n+1} - 1}{x - 1}$ \\ \hline
    $\displaystyle \sum_{i=1,2,4,\ldots,n} i$ ($n$ en potens af 2) & $\displaystyle 2n - 1$ \\ \hline
    $\displaystyle \sum_{k=0}^\infty x^k, \, |x| < 1$ & $\displaystyle \frac{1}{1 - x}$ \\ \hline
    $\displaystyle \sum_{k=1}^\infty kx^{k-1}, \, |x| < 1$ & $\displaystyle \frac{1}{(1 - x)^2}$ \\ \hline
    $\displaystyle \sum_{k=1}^n \frac{1}{k}$ & $\displaystyle \ln n + O(1)$ \\ \hline
    $\displaystyle \sum_{k=1}^n (a + bk)$ & $\displaystyle a n + b \frac{n(n+1)}{2}$ \\ \hline  
    \end{tabular}
    \caption{Nogle nyttige sumformler}
\end{table}


\section*{Sumnotation}

Betragt en følge af tal \( a_1, a_2, \ldots, a_n \), hvor \( n \) er et ikke-negativt heltal. Summen af denne følge, \( a_1 + a_2 + \dots + a_n \), kan skrives med notationen \( \sum_{k=1}^n a_k \). Når \( n = 0 \), defineres summen til at have værdien 0. Resultatet af en endelig sum er altid veldefineret, og den rækkefølge, leddene lægges sammen i, påvirker ikke slutværdien.

For en uendelig følge \( a_1, a_2, \ldots \) af tal skriver vi den uendelige sum \( a_1 + a_2 + \dots \) med notationen \( \sum_{k=1}^{\infty} a_k \), som svarer til \( \lim_{n \rightarrow \infty} \sum_{k=1}^n a_k \). Hvis denne grænseværdi eksisterer, siges rækken at være konvergent; ellers er den divergent. I modsætning til endelige summer kan leddene i en konvergent række ikke nødvendigvis byttes om uden at påvirke resultatet. I en absolut konvergent række -- dvs. en række, hvor også \( \sum_{k=1}^{\infty} |a_k| \) konvergerer -- kan leddene dog ombyttes uden at ændre summen.

\[
    \sum_{i=i_0}^n c a_i=c \sum_{i=i_0}^n a_i
\]

hvor $c$ er et vilkårligt tal. Vi kan altså sætte konstanter uden for en sum.

På samme måde kan vi dele en sum op i to summer:

\[
\sum_{i=i_0}^n\left(a_i \pm b_i\right)=\sum_{i=i_0}^n a_i \pm \sum_{i=i_0}^n b_i
\]

Bemærk, at vi har ladet rækkerne starte ved $i_0$ for at markere, at de kan starte ved en hvilken som helst værdi af $i$, vi har brug for. Bemærk også, at selv om vi kan dele summer og differenser op som ovenfor, kan vi ikke gøre det samme med produkter og kvotienter. Med andre ord:

\[
    \sum_{i=i_0}^n\left(a_i b_i\right) \neq\left(\sum_{i=i_0}^n a_i\right)\left(\sum_{i=i_0}^n b_i\right) \qquad \text{og} \qquad \sum_{i=i_0}^n \frac{a_i}{b_i} \neq \frac{\sum_{i=i_0}^n a_i}{\sum_{i=i_0}^n b_i}.
\]

\section*{Summers linearitet}
For ethvert reelt tal $c$ og alle endelige følger $a_1, a_2, \ldots, a_n$ og $b_1, b_2, \ldots, b_n$ gælder

\[
    \sum_{k=1}^n\left(c a_k+b_k\right)=c \sum_{k=1}^n a_k+\sum_{k=1}^n b_k .
\]

Linearitetsegenskaben gælder også for uendelige konvergente rækker.

Med ikke-negative funktioner og konstanter i $\Theta$-grænserne, som er uniforme i $k$, gælder samme idé for summer, der indeholder asymptotisk notation. For eksempel er

\[
    \sum_{k=1}^n \Theta(f(k))=\Theta\left(\sum_{k=1}^n f(k)\right) .
\]

I denne ligning gælder $\Theta$-notationen på venstre side variablen $k$, mens den på højre side gælder $n$. Sådanne omskrivninger gælder også for uendelige konvergente rækker.

\section*{Aritmetiske rækker}
Summen


\[
        \sum_{k=1}^n k=1+2+\cdots+n \label{eq:arithmetic-series}
\]

er en aritmetisk række og har værdien


\begin{align}
    \sum_{k=1}^n k & =\frac{n(n+1)}{2} \label{eq:sum-of-natural-numbers}  \\
    & =\Theta\left(n^2\right) . \label{eq:arithmetic-series-asymptotic}
\end{align}

En generel aritmetisk række har en additiv konstant $a \geq 0$ og en konstant koefficient $b>0$ i hvert led, men har asymptotisk samme sum:

\begin{equation}
    \sum_{k=1}^n(a+b k)=\Theta\left(n^2\right) . \label{eq:general-arithmetic-series}
    \end{equation}

\section*{Summer af kvadrater og kubiktal}
Følgende formler gælder for summer af kvadrater og kubiktal:

    \begin{equation}
    \sum_{k=0}^n k^2=\frac{n(n+1)(2 n+1)}{6}, \label{eq:sum-of-squares}
    \end{equation}

    \begin{equation}
        \sum_{k=0}^n k^3=\frac{n^2(n+1)^2}{4} . \label{eq:sum-of-cubes}
        \end{equation}

        \section*{Geometriske rækker}
        For reelle $x \neq 1$ er summen
        
        \[
            \sum_{k=0}^n x^k=1+x+x^2+\cdots+x^n
        \]
        
        en geometrisk række og har værdien
        
        \begin{equation}
            \sum_{k=0}^n x^k=\frac{x^{n+1}-1}{x-1} . \label{eq:geometric-series}
        \end{equation}
        
        Den uendelige aftagende geometriske række opstår, når summen er uendelig, og $|x|<1$:
        
        \begin{equation}
            \sum_{k=0}^{\infty} x^k=\frac{1}{1-x} . \label{eq:infinite-geometric-series}
        \end{equation}
        
        Hvis vi antager, at $0^0=1$, gælder formlerne også, når $x=0$.
        
        Et \textbf{særtilfælde} af den geometriske række opstår, når leddene er potenser af 2. Antag, at $n$ selv er en potens af 2, så leddene er \(1, 2, 4, 8, \ldots, n\). I dette tilfælde er:
        \[
            \sum_{i=1,2,4,\ldots,n} i = 1 + 2 + 4 + 8 + \cdots + n.
        \]
        
        Denne række kan udledes af den generelle formel for en geometrisk række. Den generelle form for en geometrisk række er:
        \[
            S = a + ar + ar^2 + \cdots + ar^{k-1},
        \]
        hvor:
        \begin{itemize}
            \item \(a\) er det første led (\(a=1\) i dette tilfælde),
            \item \(r\) er kvotienten (\(r=2\) her),
            \item \(k\) er antallet af led i rækken.
        \end{itemize}
        Antallet af led, \(k\), er bestemt af, hvor mange gange følgen \(1, 2, 4, \ldots, n\) fordobles, indtil den når \(n\). Da \(n = 2^{k-1}\), har vi \(k = \log_2(n) + 1\).
        
        Med formlen for den geometriske række \eqref{eq:geometric-series} får vi:
        \[
            \sum_{i=1,2,4,\ldots,n} i = \frac{2^k - 1}{2 - 1} = 2^k - 1.
        \]
        
        Ved at indsætte \(k = \log_2(n) + 1\) får vi:
        \[
            \sum_{i=1,2,4,\ldots,n} i = 2^{\log_2(n)+1} - 1 = 2n - 1.
        \]
        
        Summen af potenserne af 2 op til \(n\) er altså:
        \[
            \sum_{i=1,2,4,\ldots,n} i = 2n - 1.
        \]
        

\section*{Harmoniske rækker}
For positive heltal $n$ er det $n$'te harmoniske tal

\begin{align}
    H_n & = 1 + \frac{1}{2} + \frac{1}{3} + \frac{1}{4} + \cdots + \frac{1}{n} \nonumber \\
        & = \sum_{k=1}^n \frac{1}{k} \label{eq:harmonic_sum} \\
        & = \ln n + O(1) . \label{eq:approximation}
\end{align}

\section*{Integration og differentiation}
Ved at integrere eller differentiere formlerne ovenfor får man yderligere formler. Hvis man for eksempel differentierer begge sider af den uendelige geometriske række (\ref{eq:infinite-geometric-series}) og ganger med $x$, får man


\begin{equation}
    \begin{aligned}
    & \sum_{k=0}^{\infty} k x^k=\frac{x}{(1-x)^2} \\ 
    & \text { for }|x|<1 . \label{eq:diff-geometric-series}
    \end{aligned}
\end{equation}

\section*{Teleskoprækker}

For enhver følge $a_0, a_1, \ldots, a_n$ gælder

\begin{equation}
    \sum_{k=1}^n\left(a_k-a_{k-1}\right)=a_n-a_0, \label{eq:telescoping-series}
\end{equation}

da hvert af leddene $a_1, a_2, \ldots, a_{n-1}$ lægges til præcis én gang og trækkes fra præcis én gang. Vi siger, at summen teleskoperer. Tilsvarende gælder

\[
    \sum_{k=0}^{n-1}\left(a_k-a_{k+1}\right)=a_0-a_n .
\]

Som eksempel på en teleskopsum kan vi se på rækken

\[
\sum_{k=1}^{n-1} \frac{1}{k(k+1)}
\]

Når hvert led omskrives til

\[
\frac{1}{k(k+1)}=\frac{1}{k}-\frac{1}{k+1},
\]

får vi

\begin{align}
\sum_{k=1}^{n-1} \frac{1}{k(k+1)} & =\sum_{k=1}^{n-1}\left(\frac{1}{k}-\frac{1}{k+1}\right) \nonumber \\
& =1-\frac{1}{n} . \label{eq:telescoping-example}
\end{align}

\section*{Omindeksering af summer}
En række kan undertiden forenkles ved at ændre dens indeks, ofte ved at vende summationsrækkefølgen om. Betragt rækken $\sum_{k=0}^n a_{n-k}$. Da leddene i denne sum er $a_n, a_{n-1}, \ldots, a_0$, kan vi vende indeksernes rækkefølge om ved at sætte $j=n-k$ og omskrive summen til

$$
\sum_{k=0}^n a_{n-k}=\sum_{j=0}^n a_j
$$


Generelt er det værd at overveje en omindeksering, hvis summationsindekset optræder med minustegn i summens led.

Betragt som eksempel summen

\begin{equation}
    \sum_{k=1}^n \frac{1}{n-k+1} . \label{eq:reindexing-example}
\end{equation}


Indekset $k$ optræder med negativt fortegn i $\dfrac{1}{(n-k+1)}$. Vi kan da også forenkle denne sum, denne gang ved at sætte $j=n-k+1$, hvilket giver

\begin{equation}
    \sum_{k=1}^n \frac{1}{n-k+1}=\sum_{j=1}^n \frac{1}{j}. \label{eq:reindexing-example-solution}
\end{equation}

som netop er den harmoniske række (\ref{eq:harmonic_sum}).

\section*{Produkter}
Det endelige produkt $a_1 a_2 \ldots a_n$ kan skrives som

\[
\prod_{k=1}^n a_k .
\]


Hvis $n=0$, defineres produktets værdi til at være 1 . Man kan omdanne en formel med et produkt til en formel med en sum ved hjælp af identiteten

\[
\log \left(\prod_{k=1}^n a_k\right)=\sum_{k=1}^n \log a_k, \qquad a_k>0.
\]

\section*{Oversigtsark til tidskompleksitetsanalyse}

Dette afsnit giver de vigtigste formler og genveje til asymptotisk analyse og vurdering af algoritmer. De supplerer de detaljerede summationsmetoder i dette bilag.

\begin{itemize}
    \item \textbf{Logaritmisk tilnærmelse af fakultet (Stirlings approksimation):} 
    \[
    \ln(n!) = n\ln n-n+\mathcal{O}(\ln n).
    \]
 
    \item \textbf{Reglen for logaritmisk eksponent:} 
    \[
    2^{\log_2 n} = n.
    \]

    \item \textbf{Grundtal som potens med logaritmisk eksponent:} 
    Hvis grundtallet er \( a^b \), og eksponenten indeholder \( \log_a \), så gælder:
    \[
    (a^b)^{\log_a(n^c)} = n^{bc}, \qquad a>0,\ a\neq1,\ n>0.
    \]
    
    \textbf{Eksempel:} \( 4^{\log_2(n^2)} = n^4 \) (da \( 4 = 2^2 \)).
    
    \item \textbf{Reglen for logaritmen til en potens:} 
    \[
    \log(n^k) = k \log n, \quad \text{for } k > 0.
    \]
    
    \item \textbf{Logaritmisk vækst for potenser af 2:} 
    For \( n = 2^k \) gælder
    \[
    k = \log_2 n.
    \]
    
    \item \textbf{Forenkling af logaritmiske udtryk:} 
    \[
    \log_a n = \frac{\log_b n}{\log_b a}, \quad \text{for } a,b,n>0 \text{ med } a\neq1 \text{ og } b\neq1.
    \]
    
    \item \textbf{Almindelige asymptotiske sammenligninger:} 
    \[
    n^c \ll c^n \quad \text{for enhver konstant } c > 1.
    \]
    
    \item \textbf{Fakultetets asymptotik (udvidet Stirlings approksimation):} 
    Ifølge Stirlings formel er væksten af \( n! \) tilnærmelsesvis:
    \[
    n! \sim \sqrt{2\pi n}\left(\frac{n}{e}\right)^n,
    \]
    
    \item \textbf{Sum af potenser af 2:} 
    Hvis $n$ er en potens af 2, og leddene er \(1,2,4,8,\ldots,n\), er summen:
    \[
    \sum_{i=1,2,4,\ldots,n} i = 2n - 1.
    \]

\end{itemize}
